
This work proposes a checkpoint-gradient methodology for measuring contamination-inflation correlation in language model benchmarks. The approach uses model checkpoints across training as a natural contamination gradient and capability detrending via out-of-distribution perplexity to isolate contamination effects.

Proof-of-concept validation with simulated data demonstrates that the analysis pipeline can detect contamination-inflation correlations (Spearman r = 0.326, p = 0.003). Per-benchmark analysis shows statistically significant effects for MMLU (r = 0.55) and ARC-Challenge (r = 0.53), with non-significant results for HellaSwag and WinoGrande. Separate n-gram overlap analysis confirms that the detection mechanism functions correctly, identifying individual MMLU items with up to 17.4\% overlap in a Pile subset.

The primary limitation is that current results are from simulated validation only. Full experimental validation requires running actual Pythia checkpoint evaluations with a complete Pile n-gram index, which requires approximately 144 GPU-hours.

If confirmed by real experiments, these findings would suggest that contamination-aware benchmark evaluation is feasible. The checkpoint-gradient and capability detrending methodologies enable contamination-performance studies without requiring contamination-free baseline models.

